The tensor with the dual standard representation carries an invariant
Hermitian form of mixed signature. To estimate the norms of its
projections we first separate the diagonal and off-diagonal parts of
the split Casimir. Throughout this calculation put
\[
  \mu=\lambda+\nu,\quad
  I=\{1,\ldots,k-1\},\quad I'=\{k,\ldots,n\},\quad |I'|=m .
\]

Write \(E_{ab}\) for the complex matrix units acting on the
discrete series, and \(\varepsilon_a\) for their coordinate signs.
The compact roots join coordinates of the same sign. Let
\(\rho_c,\rho_n\) be the half sums of the positive compact and
noncompact roots, respectively, in the order defined by \(\mu\).
Thus the lowest compact type has highest weight \(L\), and its
compact Casimir eigenvalue \(K\) and the full eigenvalue \(c\) are
\begin{equation}\label{eq:paired-lowest-casimirs}
 \rho=\rho_c+\rho_n,\quad L=\mu+\rho_n-\rho_c,\quad
 K=\|\mu+\rho_n\|_2^2-\|\rho_c\|_2^2,\quad
 c=\|\mu\|_2^2-\|\rho\|_2^2 .
\end{equation}
Say that a vector has level at most \(j\) if its compact types
have compact Casimir at most \((\sqrt K+4nj)^2\), and every
diagonal weight \(\omega\) in that vector satisfies
\(\|\omega-L\|_1\le j\). The highest vector has level zero.

\begin{auxlemma}[Matrix units on the permitted compact types]\label{lem:paired-lowest-casimir}
For the weights of Lemma~\ref{lem:marked-paired-moment-range}, at
any level \(j\le4n^2+2\),
\[
  n\sum_{a\ne b}\|E_{ab}w\|^2\le (B_v/16)^2\|w\|^2,
  \qquad \sum_{a,b}\|E_{ab}w\|^2\le B_0^4\|w\|^2 .
\]
\end{auxlemma}
\begin{proof}
The adjoint relation \(E_{ab}^*=\varepsilon_a\varepsilon_b E_{ba}\)
expresses the sum of the squares as twice the compact Casimir minus
the full Casimir. For each occurring diagonal weight one has
\[
 \sum_a\omega_a^2\ge \sum_a L_a^2-2\|L\|_\infty j
 \ge\sum_a L_a^2-8B_vj .
\]
Indeed \(|\mu_a|\le 3B_v\) and
\(|L_a-\mu_a|\le n\) by the root counts.

Orthogonality of diagonal weight spaces now bounds the off-diagonal
square sum by \(D_j\|w\|^2\), where
\[
 D_j=2(\sqrt K+4nj)^2-c-\sum_a L_a^2+8B_vj .
\]

The terms of order \(B_v^2\) cancel. Coordinatewise the unshifted
summand of \(2K-c-\sum L_a^2\) is, with
\(U=(\rho_n)_a,\ T=(\rho_c)_a\),
\[
       2\mu_a(U+T)+(U+T)^2-4T^2+\rho_a^2 .
\]
The root counts give \(|U|,|T|,|U+T|,|\rho_a|\le n\).
By \eqref{eq:paired-lowest-casimirs}, \(\sqrt K\le4nB_v\).
Using \(j\le5n^2\), we obtain
\begin{align*}
 D_j &\le n(12nB_v+4n^2)+64n^2j B_v+32n^2j^2+8B_vj\\
     &\le 500n^4B_v \le \frac{B_v^2}{256n}.
\end{align*}

The last two inequalities use \(10^6 n^7\le B_v\).
The unshifted coordinates have magnitude at most \(2B_v\);
the initial shift costs at most \(B_v/100\), and the lowest-type
shift and the level cost at most \(n+j\). Hence
\[
 \max_a |\omega_a|\le 2B_v+B_v/100+n+j\le3B_v,\qquad
 9nB_v^2+B_v^2/(256n)\le B_0^4
\]
since \(B_0\ge B_v\ge10^6n^7\). Adding the diagonal squares proves
the second inequality.
\end{proof}

\begin{auxlemma}[Sum of the squared dual Pieri projections]\label{lem:paired-spectral-projections}
In the dual standard tensor, write \(R_s\) for the projection onto
the unitary summand with parameter \(\mu-e_s\), endowed with the norm
induced up to sign by the invariant Hermitian form. For any coordinate
vector \(e_c\) and any \(w\) of level \(j\le4n^2\),
\[
        \sum_s\|R_s(e_c\otimes w)\|^2
                       \le (6B_v)^{2m}\|w\|^2 .
\]
\end{auxlemma}
\begin{proof}
Write \(x=e_c\otimes w\) and let \(p_s\) project back into the tensor
product. In the positive direct-sum norm, its invariant form and
split Casimir are
\[
 [y,y]=\sum_a\varepsilon_a\|y_a\|^2,\qquad
 (Sy)_a=-\sum_b E_{ab}y_b,\qquad
 Sp_s=s_s p_s,\quad s_s=-\mu_s+(1-n)/2 .
\]
These eigenspaces and their signed isometries are the dual Pieri
decomposition used in Lemma~\ref{lem:weighted-fitted-translate}.
The distinct eigenvalues differ by integers; the eigenspaces are
orthogonal for \([\,,\,]\) since \(S\) is self-adjoint for this form.

Work on the finite-dimensional span \(V_x=\sum_s\C p_s x\).
The coordinates of elements of \(V_x\) have level at most \(j+2\).
In fact the coordinate injection, equivariant projection and
coordinate extraction have combined weight change \(e_a-e_c\).

For the compact type, the coordinate maps factor through tensoring
by the standard module and its dual. In one such tensor step the
compact infinitesimal action is \(X\otimes 1+1\otimes X\).
Using a compact orthonormal basis \((X_i)\) gives
\[
 \left(\sum_i\|X_i y\|^2\right)^{1/2}
 \le (C+\sqrt n)\|y\|
\]
on a tensor of a type of Casimir at most \(C^2\) with a standard
or dual standard module. Indeed each input Casimir is bounded there
by \(C^2\) and \(n\), and the triangle inequality applies in the
direct sum over \(i\). Compact equivariance transports this bound
to the image compact types. 
In the abstract shifted summand, a weight \(\omega-e_c\) is
measured from \(L-e_s\), so
\[
 \|(\omega-e_c)-(L-e_s)\|_1\le j+2 .
\]

The corresponding \(\sqrt K\) changes by at most \(\sqrt n\). To see the latter estimate, the two compact
highest weights differ by a standard weight; each type occurs
in the tensor of the other with the standard or its dual by
the compact Pieri rule. Thus the bound \(4n(j+2)\) in the definition
of the new level allows both tensor operations. The same argument
applies to one matrix unit, since their compact module is a summand
of the standard tensor its dual.

Write \(S=D+O\) with \((Dy)_a=-E_{aa}y_a\). The two intervals
containing the spectrum of the individual diagonal entries on these
weight spaces, and also the numbers \(s_a\), are
\[
  [-5B_v/32,B_v/32]\quad(a\in I),\qquad
  [31B_v/32,65B_v/32]\quad(a\in I').
\]
Indeed the change from \(-\lambda_a\) to a diagonal eigenvalue is
at most \(B_v/100+n+j+2\le B_v/32\).
By Cauchy--Schwarz in each row and Lemma~\ref{lem:paired-lowest-casimir},
\[
 \|Oy\|^2\le n\sum_b\sum_{a\ne b}\|E_{ab}y_b\|^2
            \le (B_v/16)^2\|y\|^2,\qquad \|Sy\|\le3B_v\|y\|
 \quad(y\in V_x).
\]

On the circle \(|\zeta+B_v/16|=9B_v/16\), both intervals have
distance at least \(15B_v/32\) from \(\zeta\). Subtracting the bound
for \(O\) and inverting on \(V_x\) gives
\[
 \|(\zeta-S)^{-1}\|\le \frac{32}{13B_v},\qquad
 P=\sum_{s\in I}p_s=\frac{1}{2\pi i}\oint(\zeta-S)^{-1}\,d\zeta,
 \qquad \|P\|\le18/13 .
\]
For \(s\in I'\) interpolate only on the remaining \(m\) eigenvalues.
Their separation is at least one, and \(\|1-P\|<3\), so
\[
 p_s=\prod_{t\in I',\ t\ne s}\frac{S-s_t}{s_s-s_t}(1-P),
 \qquad \|p_s\|\le 3(6B_v)^{m-1}.
\]

For \(s\in I\) and \(y=p_s x\), the equation
\((D-s_s)y=-Oy\) on the coordinates \(I'\) has diagonal gap
at least \(30B_v/32\). Thus
\[
  \sum_{a\in I'}\|y_a\|^2\le \|y\|^2/225,\qquad
  [y,y]\ge (1-2/225)\|y\|^2 ,
\]
because every coordinate in \(I\) has positive sign.
The positive summands therefore add without cancellation in
\([Px,Px]\), whereas each remaining signed norm is bounded by its
positive direct-sum norm. It follows that
\[
 \begin{aligned}
 \sum_s\|R_sx\|^2
 &= [Px,Px]+\sum_{s\in I'}|[p_s x,p_s x]|\\
 &\le\big((18/13)^2+9m(6B_v)^{2(m-1)}\big)\|x\|^2
 \le(6B_v)^{2m}\|w\|^2 .
 \end{aligned}
\]
Here the final comparison uses \(m\le n\le B_v\).
\end{proof}
